\chapter{Getting Started}
\label{ch:getting-started}

\begin{quote}
\textit{``The journey of a thousand miles begins with a single step.''}\\
---Lao Tzu
\end{quote}

\vspace{1em}

This chapter walks you through installing the Novus compiler, writing your first program, and understanding the compilation pipeline. By the end, you'll have a working executable running on your Amiga (or emulator).

% ============================================================================
\section{Quick Start}
\label{sec:quick-start}
% ============================================================================

If you already have Novus installed, here's the fastest path to ``Hello, Amiga!'':

\begin{enumerate}
\item Create \texttt{hello.novus}:
\begin{lstlisting}
from std::io::file import println

pub fn main() -> i32 {
    println("Hello, Amiga!")
    return 0
}
\end{lstlisting}

\item Compile:
\begin{lstlisting}[language=bash]
novus compile hello.novus -o hello
\end{lstlisting}

\item Copy to your Amiga and run:
\begin{lstlisting}[language=bash]
1.Workbench:> hello
Hello, Amiga!
\end{lstlisting}
\end{enumerate}

That's it! The rest of this chapter explains the prerequisites, installation, and toolchain in detail.

% ============================================================================
\section{Prerequisites}
\label{sec:prerequisites}
% ============================================================================

Before installing Novus, ensure you have:

\begin{itemize}
    \item \textbf{.NET 9.0 or later} --- The Novus compiler is written in C\# and requires the .NET runtime. Download from \texttt{https://dotnet.microsoft.com}

    \item \textbf{VBCC toolchain} --- Novus uses VBCC's assembler (\texttt{vasmm68k\_mot}) and linker (\texttt{vlink}) to produce Amiga executables. The toolchain is vendored in the repository under \texttt{vendor/vbcc}, or you can specify a custom installation via the \texttt{--vbcc-path} option.

    \item \textbf{An Amiga or emulator} --- To run your compiled programs. FS-UAE, WinUAE, or real hardware all work. A 68020+ configuration with AmigaOS 2.0 or later is required.
\end{itemize}

% ============================================================================
\section{Installation}
\label{sec:installation}
% ============================================================================

\subsection{Platform-Specific Instructions}

\subsubsection{macOS}

\begin{enumerate}
\item Install .NET SDK using Homebrew:
\begin{lstlisting}[language=bash]
brew install dotnet
\end{lstlisting}

\item Clone and build Novus:
\begin{lstlisting}[language=bash]
git clone https://github.com/barryw/Novus.git
cd Novus
dotnet build
\end{lstlisting}

\item Add to your PATH (optional):
\begin{lstlisting}[language=bash]
export PATH="$PATH:$(pwd)/Novus/bin/Debug/net9.0"
\end{lstlisting}
\end{enumerate}

\subsubsection{Linux}

\begin{enumerate}
\item Install .NET SDK:
\begin{lstlisting}[language=bash]
# Ubuntu/Debian
sudo apt install dotnet-sdk-9.0

# Fedora
sudo dnf install dotnet-sdk-9.0
\end{lstlisting}

\item Clone and build:
\begin{lstlisting}[language=bash]
git clone https://github.com/barryw/Novus.git
cd Novus
dotnet build
\end{lstlisting}
\end{enumerate}

\subsubsection{Windows}

\begin{enumerate}
\item Download and install .NET SDK from \texttt{https://dotnet.microsoft.com}

\item Clone the repository (using Git Bash, PowerShell, or WSL):
\begin{lstlisting}[language=bash]
git clone https://github.com/barryw/Novus.git
cd Novus
dotnet build
\end{lstlisting}
\end{enumerate}

\textbf{WSL Note:} If you prefer Linux tools, Windows Subsystem for Linux works well. Install the Linux .NET SDK inside WSL and compile there.

\subsection{Verifying Installation}

After building, verify everything works:

\begin{lstlisting}[language=bash]
# Check the Novus compiler
./Novus/bin/Debug/net9.0/novus --version

# Check VBCC tools (vendored)
./vendor/vbcc/bin/vasm --version 2>&1 | head -1
./vendor/vbcc/bin/vlink --version 2>&1 | head -1
\end{lstlisting}

You should see version information for each tool.

\subsection{Troubleshooting}

\textbf{Problem:} \texttt{dotnet: command not found}\\
\textbf{Solution:} .NET SDK is not installed or not in PATH. Reinstall and ensure the install location is in your PATH.

\textbf{Problem:} \texttt{VBCC tools not found}\\
\textbf{Solution:} The vendored VBCC is under \texttt{vendor/vbcc}. If using a custom VBCC installation, pass \texttt{--vbcc-path /path/to/vbcc} to the compiler.

\textbf{Problem:} Build fails with missing dependencies\\
\textbf{Solution:} Run \texttt{dotnet restore} to fetch NuGet packages.

\begin{comingfromc}
Setting up a C cross-compiler for Amiga traditionally involves:
\begin{enumerate}
\item Downloading VBCC, vasm, and vlink separately
\item Setting \texttt{VBCC} environment variable
\item Creating \texttt{vc.config} with include/library paths
\item Finding and installing the AmigaOS NDK
\item Setting up include paths for system headers
\end{enumerate}

Novus bundles everything in the repository---no environment variables or configuration files needed. Just \texttt{dotnet build} and compile.
\end{comingfromc}

% ============================================================================
\section{Your First Program}
\label{sec:first-program}
% ============================================================================

Create a file named \texttt{hello.novus}:

\begin{lstlisting}
from std::io::file import println

pub fn main() -> i32 {
    println("Hello, Amiga!")
    return 0
}
\end{lstlisting}

\subsection{Line-by-Line Explanation}

\textbf{Line 1: \texttt{from std::io::file import println}}

This imports the \texttt{println} function from the standard library's I/O module. Unlike some languages where printing is built-in, Novus requires explicit imports. This keeps the language core small and makes dependencies visible.

\texttt{std::io::file} is the module path:
\begin{itemize}
    \item \texttt{std} --- the standard library
    \item \texttt{io} --- the I/O subsystem
    \item \texttt{file} --- file and console operations
\end{itemize}

\textbf{Line 3: \texttt{pub fn main() -> i32}}

This declares the program entry point. Every Novus program must have exactly one \texttt{main} function that:
\begin{itemize}
    \item Is marked \texttt{pub} (public) so the linker can find it
    \item Takes no parameters
    \item Returns \texttt{i32} --- the exit code passed back to AmigaOS
\end{itemize}

\textbf{Line 4: \texttt{println("Hello, Amiga!")}}

Calls the imported \texttt{println} function with a string literal. String literals in Novus are null-terminated C-style strings (\texttt{*u8}).

\textbf{Line 5: \texttt{return 0}}

Returns the exit code. Zero means success; non-zero indicates an error. AmigaOS and the Shell use this to detect failures.

\subsection{Compiling}

Compile to an Amiga executable:

\begin{lstlisting}[language=bash]
novus compile hello.novus -o hello
\end{lstlisting}

This produces an executable named \texttt{hello} in Amiga HUNK format. Without the \texttt{-o} flag, the output defaults to \texttt{a.out}.

\subsection{What Happens During Compilation}

When you compile, Novus performs these steps:

\begin{enumerate}
\item \textbf{Parse} --- Read \texttt{hello.novus} and build an AST
\item \textbf{Analyze} --- Check types, resolve names, build IR
\item \textbf{Import} --- Load \texttt{std::io::file} and its dependencies
\item \textbf{Generate} --- Emit C99 code for your program and used stdlib functions
\item \textbf{Compile} --- VBCC compiles C to 68k object files
\item \textbf{Link} --- vlink combines objects into a HUNK executable
\end{enumerate}

Use \texttt{-v} (verbose) to watch this process:

\begin{lstlisting}[language=bash]
novus compile hello.novus -o hello -v
\end{lstlisting}

\subsection{Running on the Amiga}

Transfer the executable to your Amiga and run it from the Shell:

\begin{lstlisting}[language=bash]
1.Workbench:> hello
Hello, Amiga!
\end{lstlisting}

Congratulations---you've run your first Novus program!

% ============================================================================
\section{A More Complete Example}
\label{sec:complete-example}
% ============================================================================

Let's look at a program that does more than print a message. This example queries and displays system memory information:

\begin{lstlisting}
// memory_info.novus - Display memory statistics
from std::io::file import print, println
from std::ffi::exec import AvailMem
from std::ffi::amiga_consts import MEMF_CHIP, MEMF_FAST, MEMF_TOTAL

pub fn main() -> i32 {
    println("Memory Information")
    println("==================")

    // Query chip RAM (DMA-accessible memory)
    let chip_mem = unsafe { AvailMem(MEMF_CHIP | MEMF_TOTAL) }

    // Query fast RAM (CPU-only memory)
    let fast_mem = unsafe { AvailMem(MEMF_FAST | MEMF_TOTAL) }

    // Display results
    print("Chip RAM: ")
    print_kb(chip_mem)
    println(" KB")

    print("Fast RAM: ")
    print_kb(fast_mem)
    println(" KB")

    print("Total:    ")
    print_kb(chip_mem + fast_mem)
    println(" KB")

    return 0
}

fn print_kb(bytes: u32) {
    let kb = bytes / 1024u32
    print_number(kb)
}

fn print_number(n: u32) {
    if n == 0 { return }
    if n >= 10 { print_number(n / 10u32) }
    let digit: u8 = (u8)(n % 10u32) + 48u8
    var buf = [0u8; 2]
    buf[0] = digit
    print(&buf[0])
}
\end{lstlisting}

This example demonstrates:
\begin{itemize}
    \item \textbf{Multiple imports} from different modules
    \item \textbf{Unsafe FFI calls} to AmigaOS via \texttt{AvailMem}
    \item \textbf{Constants} like \texttt{MEMF\_CHIP} from the FFI layer
    \item \textbf{Helper functions} for formatting output
    \item \textbf{Arithmetic and type casts}
\end{itemize}

\begin{comingfromc}
The equivalent C code would look similar, but you'd need to:
\begin{itemize}
    \item Include \texttt{<proto/exec.h>} or set up function prototypes
    \item Link with \texttt{amiga.lib} or use direct library calls
    \item Deal with library bases (\texttt{SysBase})
\end{itemize}

In Novus, the FFI layer handles library opening/closing automatically. Just import and call.
\end{comingfromc}

% ============================================================================
\section{The Compilation Pipeline}
\label{sec:compilation-pipeline}
% ============================================================================

Understanding the compilation pipeline helps with debugging and optimization. Here's what happens when you compile:

% TikZ diagram of compilation pipeline
\begin{center}
\begin{tikzpicture}[
    node distance=0.8cm,
    box/.style={rectangle, draw, rounded corners, minimum width=2.2cm, minimum height=0.7cm, align=center, font=\small},
    arrow/.style={->, >=stealth, thick}
]
    \node[box] (source) {.novus\\source};
    \node[box, right=of source] (ir) {IR};
    \node[box, right=of ir] (c) {.c\\C99};
    \node[box, right=of c] (obj) {.o\\objects};
    \node[box, right=of obj] (exe) {HUNK\\executable};

    \draw[arrow] (source) -- node[above, font=\scriptsize] {parse} (ir);
    \draw[arrow] (ir) -- node[above, font=\scriptsize] {codegen} (c);
    \draw[arrow] (c) -- node[above, font=\scriptsize] {vbcc} (obj);
    \draw[arrow] (obj) -- node[above, font=\scriptsize] {vlink} (exe);
\end{tikzpicture}
\end{center}

\subsection{Stage 1: Parsing and Analysis}

The Novus compiler reads your source file, parses it into an Abstract Syntax Tree (AST), then performs semantic analysis:
\begin{itemize}
    \item Type checking
    \item Name resolution
    \item Trait resolution
    \item Building the Intermediate Representation (IR)
\end{itemize}

Use \texttt{--emit-ir} to see the IR:

\begin{lstlisting}[language=bash]
novus compile hello.novus --emit-ir
\end{lstlisting}

\subsection{Stage 2: C Code Generation}

The IR is lowered to C99 code. Each function becomes a C function; each type becomes a C struct or typedef. The generated C uses VBCC-specific extensions for 68k calling conventions.

\subsection{Stage 3: Object File Compilation}

VBCC's C compiler (\texttt{vbccm68k}) compiles the C code to 68k object files. Each function is compiled separately for incremental builds.

Use \texttt{--emit-asm} to see the generated assembly:

\begin{lstlisting}[language=bash]
novus compile hello.novus --emit-asm
\end{lstlisting}

This is useful for understanding performance characteristics or debugging codegen issues.

\subsection{Stage 4: Linking}

\texttt{vlink} combines all object files---your code, the runtime, and stdlib---into a final executable in Amiga HUNK format. Dead code elimination removes unused functions to keep binaries small.

% ============================================================================
\section{Development Workflow}
\label{sec:dev-workflow}
% ============================================================================

\subsection{Fast Feedback with novus check}

For quick syntax and type checking without full compilation:

\begin{lstlisting}[language=bash]
novus check myfile.novus
\end{lstlisting}

This skips code generation and linking, giving faster feedback during development.

\subsection{Code Formatting}

Keep your code consistent with \texttt{novus fmt}:

\begin{lstlisting}[language=bash]
novus fmt myfile.novus
\end{lstlisting}

This reformats the file in place according to the Novus style guide.

\subsection{Editor Support}

Novus provides a Language Server Protocol (LSP) implementation for IDE integration:
\begin{itemize}
    \item \textbf{VS Code} --- Install the Novus extension from the marketplace
    \item \textbf{Other editors} --- Configure your LSP client to use \texttt{novus lsp}
\end{itemize}

Features include:
\begin{itemize}
    \item Syntax highlighting
    \item Error diagnostics as you type
    \item Go to definition
    \item Hover information
    \item Code completion
\end{itemize}

% ============================================================================
\section{Running on Real Hardware}
\label{sec:real-hardware}
% ============================================================================

The ultimate goal is running on real Amiga hardware. Here are several ways to transfer executables.

\subsection{Shared Folders (Emulator)}

If using FS-UAE, WinUAE, or Amiberry, configure a shared folder:

\begin{enumerate}
\item Create a directory on your host system (e.g., \texttt{\textasciitilde/amiga\_share})
\item Configure the emulator to mount it as a device (e.g., \texttt{DEV:})
\item Compile directly to the shared folder:
\begin{lstlisting}[language=bash]
novus compile hello.novus -o ~/amiga_share/hello
\end{lstlisting}
\item Access from the Amiga: \texttt{DEV:hello}
\end{enumerate}

\subsection{Network Transfer}

With networked Amigas (via Plipbox, A2065, or similar):

\begin{lstlisting}[language=bash]
# Using scp (requires AmiSSH or similar on the Amiga)
scp hello amiga.local:Work/

# Using ftp
ftp amiga.local
put hello
\end{lstlisting}

\subsection{CF/SD Card}

For accelerator cards with CF/SD storage:

\begin{enumerate}
\item Mount the card on your development machine
\item Copy the executable to the card
\item Unmount and return the card to the Amiga
\end{enumerate}

\subsection{Floppy Disk Images}

Create an ADF image with your executable:

\begin{lstlisting}[language=bash]
# Create a new ADF
xdftool myfiles.adf format "MyDisk" + write hello
\end{lstlisting}

Then use the ADF in your emulator or write it to a real floppy.

% ============================================================================
\section{Project Structure}
\label{sec:project-structure}
% ============================================================================

For anything beyond a single file, create a project with \texttt{novus new}:

\begin{lstlisting}[language=bash]
novus new myproject --type cli
cd myproject
\end{lstlisting}

This creates a standard project structure:

\begin{lstlisting}[language=bash]
myproject/
    project.toml        # Project configuration
    src/
        main.novus      # Entry point
\end{lstlisting}

\subsection{The project.toml File}

The \texttt{project.toml} file defines project metadata and build settings:

\begin{lstlisting}
[package]
name = "myproject"
version = "0.1.0"
authors = ["Your Name"]
description = "My first Novus project"

[build]
target_cpu = "68020"
fpu = "auto"
optimization_level = 0
output = "build"
\end{lstlisting}

\subsection{Build Settings Reference}

\begin{center}
\begin{tabular}{llp{6cm}}
\textbf{Setting} & \textbf{Default} & \textbf{Description} \\
\hline
\texttt{target\_cpu} & \texttt{auto} & Target CPU: 68020, 68030, 68040, 68060, or auto \\
\texttt{fpu} & \texttt{auto} & FPU mode: soft, 68881, 68040, or auto \\
\texttt{optimization\_level} & \texttt{0} & VBCC optimization level (0--3) \\
\texttt{output} & \texttt{build} & Output directory for compiled files \\
\texttt{stack\_size} & \texttt{4096} & Stack size in bytes \\
\end{tabular}
\end{center}

Build the project with:

\begin{lstlisting}[language=bash]
novus build
\end{lstlisting}

% ============================================================================
\section{Compiler Commands Reference}
\label{sec:commands}
% ============================================================================

\begin{description}
    \item[\texttt{novus compile <file>}] Compile a single source file to an executable
    \item[\texttt{novus build}] Build a project defined by \texttt{project.toml}
    \item[\texttt{novus check <file>}] Type-check without generating code
    \item[\texttt{novus test <file>}] Compile test runner
    \item[\texttt{novus new <name>}] Create a new project from a template
    \item[\texttt{novus fmt <file>}] Format source code
    \item[\texttt{novus clean}] Remove cached build artifacts
    \item[\texttt{novus lsp}] Start the language server
\end{description}

\subsection{Common Options}

Options for \texttt{novus compile}:

\begin{description}
    \item[\texttt{-o, --output}] Output file name (default: \texttt{a.out})
    \item[\texttt{--cpu}] Target CPU: \texttt{68020}, \texttt{68030}, \texttt{68040}, \texttt{68060}, or \texttt{auto}
    \item[\texttt{--fpu}] FPU mode: \texttt{soft}, \texttt{68881}, \texttt{68040}, or \texttt{auto}
    \item[\texttt{--release}] Enable optimizations and disable debug checks
    \item[\texttt{-O}] Optimization level (0--2)
    \item[\texttt{--emit-asm}] Emit assembly output only
    \item[\texttt{--emit-ir}] Emit IR to stdout
    \item[\texttt{-v, --verbose}] Show detailed progress
\end{description}

\subsection{Debug vs Release Builds}

By default, Novus compiles in debug mode:
\begin{itemize}
    \item Array bounds checking enabled
    \item Debug symbols included
    \item Minimal optimizations
\end{itemize}

For production, use \texttt{--release}:

\begin{lstlisting}[language=bash]
novus compile hello.novus -o hello --release
\end{lstlisting}

% ============================================================================
\section{The Standard Library}
\label{sec:stdlib}
% ============================================================================

Novus includes a comprehensive standard library:

\begin{description}
    \item[\texttt{std::core}] Fundamental types: \texttt{Option}, \texttt{Result}, core traits
    \item[\texttt{std::collections}] Data structures: \texttt{Vec}, \texttt{HashMap}, \texttt{HashSet}
    \item[\texttt{std::strings}] String handling: \texttt{String}, \texttt{Str}, \texttt{StringBuilder}
    \item[\texttt{std::memory}] Memory management, slices
    \item[\texttt{std::math}] Fixed-point arithmetic, trigonometry
    \item[\texttt{std::io}] File and console I/O
    \item[\texttt{std::ffi}] AmigaOS bindings (Exec, DOS, Graphics, Intuition)
    \item[\texttt{std::hardware}] CPU/chipset detection, memory info
\end{description}

Import with \texttt{from}:

\begin{lstlisting}
from std::collections::vec import Vec
from std::io::file import println
\end{lstlisting}

% ============================================================================
\section{Showcase: System Information}
\label{sec:showcase-sysinfo}
% ============================================================================

This program demonstrates the concepts from this chapter by querying and displaying system information:

\begin{lstlisting}
// sysinfo.novus - System Information Display
from std::core import Option
from std::io::file import print, println
from std::ffi::exec import AvailMem
from std::ffi::amiga_consts import MEMF_FAST, MEMF_CHIP
from std::ffi::amiga_consts import MEMF_TOTAL, MEMF_LARGEST

fn print_memory(label: *u8, bytes: u32) {
    print("  ")
    print(label)
    print(": ")
    let kb: u32 = bytes / 1024u32
    if kb == 0 {
        print("0")
    } else {
        print_number(kb)
    }
    println(" KB")
}

fn print_number(n: u32) {
    if n == 0 { return }
    if n >= 10 { print_number(n / 10u32) }
    let digit: u8 = (u8)(n % 10u32) + 48u8
    var buf = [0u8; 2]
    buf[0] = digit
    print(&buf[0])
}

fn show_memory_info() {
    println("")
    println("MEMORY")
    println("----------------------------------------")

    let chip_total = unsafe { AvailMem(MEMF_CHIP | MEMF_TOTAL) }
    let chip_largest = unsafe { AvailMem(MEMF_CHIP | MEMF_LARGEST) }
    print_memory("Chip RAM Total", chip_total)
    print_memory("Chip RAM Largest Block", chip_largest)

    println("")

    let fast_total = unsafe { AvailMem(MEMF_FAST | MEMF_TOTAL) }
    if fast_total > 0 {
        let fast_largest = unsafe { AvailMem(MEMF_FAST | MEMF_LARGEST) }
        print_memory("Fast RAM Total", fast_total)
        print_memory("Fast RAM Largest Block", fast_largest)
    } else {
        println("  Fast RAM: Not installed")
    }
}

pub fn main() -> i32 {
    println("============================================")
    println("         NOVUS SYSTEM INFORMATION")
    println("============================================")

    show_memory_info()

    println("")
    println("============================================")
    return 0
}
\end{lstlisting}

This program demonstrates:
\begin{itemize}
    \item \textbf{Imports} from multiple stdlib modules
    \item \textbf{Unsafe FFI calls} to \texttt{AvailMem}
    \item \textbf{Constants} from \texttt{amiga\_consts}
    \item \textbf{Helper functions} for output formatting
    \item \textbf{Control flow} with conditionals
\end{itemize}

The full source code is available in \texttt{Novus.Tests/Examples/sysinfo\_showcase.novus}.

% ============================================================================
\section{Running Tests}
\label{sec:running-tests}
% ============================================================================

Novus has built-in test support. Mark test functions with the \texttt{@test} attribute:

\begin{lstlisting}
from std::test::test import expect_eq

@test
fn test_addition() {
    expect_eq(2 + 2, 4, "basic addition")
}

@test
fn test_multiplication() {
    expect_eq(3 * 7, 21, "basic multiplication")
}
\end{lstlisting}

Compile the test runner:

\begin{lstlisting}[language=bash]
novus test mytest.novus
\end{lstlisting}

\textbf{Important:} The test command builds the test runner but does not execute it. Transfer the executable to your Amiga to run the tests.

% ============================================================================
\section{Next Steps}
\label{sec:next-steps}
% ============================================================================

You now have a working Novus installation and understand the basic workflow. The following chapters cover:

\begin{itemize}
    \item \textbf{Chapter~\ref{ch:basics}: Language Basics} --- Variables, types, and expressions
    \item \textbf{Chapter~\ref{ch:types}: Types} --- The type system in depth
    \item \textbf{Chapter~\ref{ch:memory}: Memory Management} --- Safe and explicit memory handling
    \item \textbf{Chapter~\ref{ch:control-flow}: Control Flow} --- Conditionals, loops, and pattern matching
\end{itemize}

When you're ready, turn the page and start learning the language itself.
